\subsection{Amount of tissue}
\label{sec:amount-discussion}
\subsubsection{Concentration and size cannot be separated}
The fluorescence shows only the product of
concentration and size: a small piece with a high concentration can give the same peak as a larger piece
with a lower one. Since concentrations are mixed in surgery, the area cannot be converted to
volume. The peak height does not help either, because the tumour pieces seem to be smaller
(more pieces, lower peaks, same area). This needs a second signal that depends on size but not
on the dye, such as scattered or transmitted light.

\subsubsection{Widths depend on the flow rate}
The concentration-free widths failed because the flow
rate changes between sets and even within one recording (the tube sometimes gets stuck). The
pulse-width method also assumes a light spot smaller than the piece~\cite{sharpless1976},
which is probably not true for small tumour pieces. A flow sensor and the size of the light
spot would be needed.

\subsubsection{The count depends on how tissue breaks}
The count does not depend on the
concentration and gave 11--32\,\% error, but it changes with how finely the tissue breaks and
with pieces joining in dense suctions (Section~\ref{sec:detection-discussion}).

\subsubsection{Calibration data}
The data sets have four or five recordings per group, so the errors
themselves are uncertain. Recordings of single pieces with known size and concentration, with the
flow rate, gain, dye concentration and tissue weight noted for each recording, would be needed
to test whether a calibration to volume is possible at all.
